% ====================================================================
%  现有技术检索记录（内部工作文档，非递交件）
%  由 notes.sty 排版。要改内容直接改本文件，这是唯一来源。
% ====================================================================
\documentclass[11pt,UTF8]{ctexart}
\usepackage{notes}
\notestitle{现有技术检索记录}

\begin{document}

\nH{现有技术检索记录（A1 / A2 / A3）}

\begin{nquote}
{\ntbf 来源}：本文件如实转录申请人在本轮检索中提供的结果，未由本工具独立复核。 {\ntbf 用途}：供代理人复核、供背景技术与创造性论证参考。 {\ntbf 禁止事项}：未经代理人核实，下列文献号不得直接写入说明书背景技术或答复审查意见。
\end{nquote}

\nHH{1. 与 A1（工具调用循环检测与阻断）相关}

\begin{ntable}{Y{0.3}Y{0.3}Y{0.3}}
{\ntbf 文献 / 出处} & {\ntbf 命中的特征} & {\ntbf 对 A1 的影响} \\
\nthead
arXiv 2603.10060《Tool Receipts》 & runtime 代执行工具并生成签名回执，回执同时含 \nc{input\_hash}（入参 canonical JSON 哈希）与 \nc{output\_hash}（原始输出哈希）。用途为防幻觉溯源，非循环检测 & {\ntbf 占用}"规范化 JSON 后哈希"+"哈希输出侧"两项特征。A1 不得把这两点单独作为发明点 \\
US20110280145A1（网络包环路检测） & 将地址与载荷代入哈希函数作 KEY，对该哈希值的计数字段累加，计数表示环路候选被观测次数 & {\ntbf 占用}"内容哈希 + 计数判环"。最可能的创造性驳回引证 \\
US20080235163（爬虫去重） & 内容哈希去重 & 同上，跨领域常规应用风险 \\
US20100205265A1（邮件去重） & 去重时刻意剔除时间戳后再哈希 & {\ntbf 占用}"剔除易变字段后哈希"。属跨领域，尚可争辩但不宜作为唯一支点 \\
\end{ntable}

{\ntbf 未命中（A1 幸存的差异特征）}：

\begin{enumerate}
  \item 产物内容摘要值同时充当证据身份（路径含调用标识故不可用）；
  \item 三重合取阻断条件（连续无新证据 ∧ 调用重复≥3 ∧ 结果重复≥3），现有技术均为单条件 k 次；
  \item 账本可从消息历史重建以跨越上下文压缩；
  \item 轮询类工具阈值整数倍放大。
\end{enumerate}

\nHH{2. 与 A2（大输出外置）相关 —— 已被打穿，不能独立申请}

\begin{ntable}{Y{0.46}Y{0.46}}
{\ntbf 文献 / 出处} & {\ntbf 命中的特征} \\
\nthead
WO2025221751A1《Dynamic context window management for conversational AI agents》 & prompt 预算按可用上下文窗口动态分配；完整会话历史载荷保留在 memory service，仅把预算内子集装入 prompt；预算随工具调用返回迭代调整 \\
MemGPT（虚拟上下文分页，2023 公开） & 绕不过去的 102/103 对比文件 \\
arXiv 2606.17016 TokenPilot & 占位符 + recovery tool；明确指出 recovery tool 是必需的 \\
arXiv 2604.08224（综述） & "上下文是最稀缺的共享资源、各类内容竞争有限 token 预算"写为领域常识 —— 对创造性杀伤力极大 \\
\end{ntable}

{\ntbf 结论}：A2 降级为 A1 的从属权利要求。独立申请必死于 MemGPT + WO2025221751A1。

{\ntbf 救法}：以"外置后用同一内容摘要值充当证据身份，使外置不破坏循环检测"这一桥接特征并入 A1 独权（检索无命中，且为真正的协同效应而非功能拼接）。

\nHH{3. 与 A3（互助出图网络）相关 —— 骨架被完整覆盖}

\begin{ntable}{Y{0.46}Y{0.46}}
{\ntbf 文献} & {\ntbf 命中的特征} \\
\nthead
US12373712《AI model inference and tuning peer-to-peer network system》 & P2P signal node 居间撮合 client 与分布式计算节点，交换 offer 建立连接，结果直接回传不经 signal node，校验值不符则从注册表注销节点 \\
US9992342B1《System for providing remote expertise》 & 撮合服务通知匹配专家、首个接受者胜出、SDP 建 P2P 连接或改走 relay、已匹配专家标记占用不再参与后续撮合 \\
US20220006860A1 / WO2020099924A1（Intel 系） & 去中心化算力市场、managing node 撮合、分层选节点 \\
\end{ntable}

{\ntbf 幸存差异}：中继铸造权单向归请求方（helper 无铸造权）、双边 settled 语义、借用 helper 自有第三方账号而凭据不出本机。

{\ntbf 结论}：A3 若申请，独权必须窄到上述三点，"撮合 + P2P + relay 回退"只能写为前序技术特征。

\nHH{4. 与 B5（MATLAB 文件 IPC）相关}

无直接命中"MATLAB + 文件式请求响应 + 持久会话"。但 MatlabMPI 文件式消息传递（学术文献）、MathWorks MATLAB Engine API、US9733992B1（容器间设备文件路由）、CN102508722B（文件发布订阅）构成"常规技术选择"驳回的基础。新颖性大概率有，创造性大概率被驳。

\nHH{5. 检索消除不掉的风险：抵触申请}

专利申请有 18 个月公开延迟。以 Anthropic 为申请人检索无命中，更可能是 2025 年的申请尚未到公开期，而非未申请。依《专利法》第 22 条第 2 款，在先申请、在后公开的同族申请构成{\ntbf 抵触申请}，破坏新颖性且任何检索都查不到。

该风险对 A1 尤其实在：OpenAI / Anthropic / 微软 / 字节均在做 agent runtime，循环防护是其必然遇到的问题。{\ntbf 无法规避，只能靠早申请。}

\nHH{6. 申请人自身在先公开（最高优先级风险）}

代码仓库已于 2026-09-13 公开推送（\nc{origin/0.4.69}，提交 \nc{fedd467f}），其中 \nc{crates/runtime/src/evidence\_ledger.rs}、\nc{crates/runtime/src/tool\_output\_artifact.rs} 及 \nc{docs/development-logic/evidence-ledger.md} 完整披露了本申请的技术方案。

《专利法》第 24 条的六个月宽限期{\ntbf 仅}适用于：国家紧急状态下为公共利益首次公开、中国政府主办或承认的国际展览会首次展出、规定的学术会议或技术会议首次发表、他人未经申请人同意泄露。{\ntbf 申请人自行在代码托管平台公开不属于任何一种情形，不享有宽限期。}

{\ntbf 结论}：自 2026-09-13 起，该公开即构成本申请的现有技术。递交日每推迟一天，风险只增不减。{\ntbf 以最快速度递交是唯一的缓解手段。}

\nHH{7. 排序结论}

{\ntbf A1（含 A2 从属） > A3（窄化版） > B5}

选 A1 的理由：

\begin{enumerate}
  \item 三个精确检索式返回 0 件，中文库无一件覆盖完整链路，新颖性风险最低；
  \item 技术效果落在《专利审查指南》明确认可的"提升计算机系统内部性能"一类（减少无效工具调用与网络请求、降低算力与 token 消耗），基本不会被以"智力活动的规则和方法"驳回；
  \item "账本必须能从消息历史重建，否则压缩后就忘了自己在打转"是我方独有的技术问题——他人做循环检测时上下文不会被压缩——这是最好的创造性支点；
  \item A2 在此复活：作为独权中的桥接特征，检索无命中且属真正的协同效应。
\end{enumerate}

\end{document}
